\chapter{Projections}

Suppose you stand in a flat field off a straight road and want the shortest walk
to the road. You walk straight toward it, meeting the road at a right angle.
That foot of the perpendicular is the closest point -- the \emph{orthogonal
projection}.

Nothing about that picture changes when we move from a line in
$\R^2$ to a plane in $\R^3$ or an $r$-dimensional subspace inside $\R^n$: to
find the point in a subspace closest to $\vec{b}$, drop a perpendicular.

The reason projections matter so much is that linear systems in practice are
often overdetermined. In $\mat{A}\vec{x}=\vec{b}$, when $\vec{b}$ is not in the
column space $\col(\mat{A})$, there is no exact solution. Instead of seeking an
exact solution that does not exist, we find the closest reachable vector: the
orthogonal projection of $\vec{b}$ onto $\col(\mat{A})$. This is the foundation
of least-squares approximation, linear regression, and the normal equations.

\section{Projecting onto a line}

Start with the simplest case: a single line through the origin, and a point we
want to approximate. Let the line be the span of a vector $\vec{a}$, and let
$\vec{b}$ be the point we are trying to reach. The line is
$\set{\,t\vec{a} : t \in \R\,}$ -- all scalar multiples of $\vec{a}$. We seek
the point $\vec{p}$ on this line that minimizes the distance $\norm{\vec{b}-\vec{p}}$.

\begin{intuitionbox}
The projection $\vec{p}$ is the point on the line closest to $\vec{b}$. The segment
connecting $\vec{b}$ to $\vec{p}$ meets the line at an exact right angle. This
perpendicular vector is the \emph{error} $\vec{e} = \vec{b} - \vec{p}$, and
demanding that the error be orthogonal to $\vec{a}$ uniquely determines $\vec{p}$.
\end{intuitionbox}

\begin{figure}[H]
\centering
\begin{tikzpicture}[scale=1.4,>=Stealth]
  % the line through origin (direction a)
  \draw[spacegray,thick] (-0.6,-0.3)--(3.5,1.75);
  \node[spacegray,below right] at (3.3,1.65){$\text{line} = \spn(\vec{a})$};
  % a
  \draw[->,vecblue,very thick] (0,0)--(1.1,0.55);
  \node[vecblue,above left] at (0.6,0.38){$\vec{a}$};
  % p = projection (foot of perpendicular)
  \draw[->,targetgreen,very thick] (0,0)--(2.0,1.0);
  \node[targetgreen,below right] at (2.0,0.95){$\vec{p}=\hat{x}\vec{a}$};
  % b
  \draw[->,vecred,very thick] (0,0)--(1.25,2.5);
  \node[vecred,above] at (1.25,2.55){$\vec{b}$};
  % error e = b - p, dashed, perpendicular
  \draw[quotientorange,very thick,dashed] (1.25,2.5)--(2.0,1.0);
  \node[quotientorange,above right] at (1.65,1.8){$\vec{e}=\vec{b}-\vec{p}$};
  % right angle mark at p: exact right angle with slope 0.5 and -2
  \draw[spacegray,thick] (2.0,1.0) ++(-0.16,-0.08) -- ++(-0.08,0.16) -- ++(0.16,0.08);
  \fill (2.0,1.0) circle (1.2pt);
\end{tikzpicture}
\caption{Projection onto a line. The shadow $\vec{p}$ is the point on $\spn(\vec{a})$ closest to
$\vec{b}$. The error $\vec{e}=\vec{b}-\vec{p}$ (orange, dashed) meets the line at an exact right angle.}
\label{ch18:project-line}
\end{figure}

\subsection{Finding the projection from the right angle}

Because $\vec{p}$ lies on the line, it is \emph{some} multiple of $\vec{a}$:
\[
  \vec{p} = \hat{x}\,\vec{a}
\]
for an unknown scalar $\hat{x}$ (read ``$x$-hat''; the hat will always mean ``our best estimate'').
The whole problem is to find this one number. We have one fact to pin it down: the error
$\vec{e} = \vec{b} - \vec{p} = \vec{b} - \hat{x}\vec{a}$ must be \emph{perpendicular to the line},
which means perpendicular to $\vec{a}$. Perpendicular means the dot product is zero:
\[
  \vec{a}\T(\vec{b} - \hat{x}\vec{a}) = 0.
\]
That is one linear equation in the one unknown $\hat{x}$. Expand it:
\[
  \vec{a}\T\vec{b} - \hat{x}\,\vec{a}\T\vec{a} = 0
  \qquad\Longrightarrow\qquad
  \hat{x} = \frac{\vec{a}\T\vec{b}}{\vec{a}\T\vec{a}}.
\]
And then the projection itself is
\[
  \boxed{\;\vec{p} = \hat{x}\,\vec{a} = \frac{\vec{a}\T\vec{b}}{\vec{a}\T\vec{a}}\,\vec{a}.\;}
\]

\begin{focusbox}
The projection of $\vec{b}$ onto the line $\spn(\vec{a})$ is
\[
  \vec{p} = \frac{\vec{a}\T\vec{b}}{\vec{a}\T\vec{a}}\,\vec{a},
  \qquad
  \hat{x} = \frac{\vec{a}\T\vec{b}}{\vec{a}\T\vec{a}}.
\]
The numerator $\vec{a}\T\vec{b}$ measures how much $\vec{b}$ leans along $\vec{a}$; the denominator
$\vec{a}\T\vec{a}=\norm{\vec{a}}^2$ normalizes for how long $\vec{a}$ is. Notice $\vec{p}$ does not
change if you rescale $\vec{a}$ -- the line is what matters, not the particular arrow on it.
\end{focusbox}

\begin{example}\label{ch18:line-numeric}
Project $\vec{b}=\colvec{2\\3}$ onto the line through $\vec{a}=\colvec{1\\0}$ (the $x$-axis).
\[
  \vec{a}\T\vec{b} = 1\cdot 2 + 0\cdot 3 = 2,
  \qquad
  \vec{a}\T\vec{a} = 1,
  \qquad
  \hat{x} = \tfrac{2}{1} = 2,
\]
so $\vec{p} = 2\colvec{1\\0} = \colvec{2\\0}$ and the error is
$\vec{e} = \vec{b}-\vec{p} = \colvec{0\\3}$. Of course: projecting onto the $x$-axis just keeps the
$x$-coordinate and drops the $y$-coordinate. And indeed $\vec{e}=\colvec{0\\3}$ is perpendicular to
$\vec{a}=\colvec{1\\0}$, since their dot product is $0$. The formula reproduces what you already
knew, which is exactly what a good formula should do on an easy case.
\end{example}

\begin{example}\label{ch18:line-diagonal}
Now a case you could not eyeball. Project $\vec{b}=\colvec{1\\4}$ onto the line through
$\vec{a}=\colvec{2\\1}$.
\[
  \vec{a}\T\vec{b} = 2\cdot1 + 1\cdot4 = 6,
  \qquad
  \vec{a}\T\vec{a} = 2^2 + 1^2 = 5,
  \qquad
  \hat{x} = \tfrac{6}{5}.
\]
So
\[
  \vec{p} = \tfrac{6}{5}\colvec{2\\1} = \colvec{12/5\\6/5},
  \qquad
  \vec{e} = \vec{b}-\vec{p} = \colvec{1-12/5\\ 4-6/5} = \colvec{-7/5\\14/5}.
\]
Check perpendicularity: $\vec{a}\T\vec{e} = 2(-7/5)+1(14/5) = -14/5+14/5 = 0$. The error really is
orthogonal to the line. This check -- ``dot the error with $\vec{a}$ and get zero'' -- is the single
most reliable way to catch an arithmetic slip in any projection problem.
\end{example}

\begin{walkbox}
Take the formula $\vec{p} = \dfrac{\vec{a}\T\vec{b}}{\vec{a}\T\vec{a}}\,\vec{a}$ on a walk and ask:
what happens if $\vec{b}$ is already \emph{on} the line, say $\vec{b}=5\vec{a}$? What should the
shadow be, and does the formula give it? And what if $\vec{b}$ is perpendicular to $\vec{a}$ -- where
does the shadow land then? Two sanity checks that a projection must pass.
\end{walkbox}

\subsection{Packaging the projection as a matrix}

Look again at $\vec{p} = \dfrac{\vec{a}\T\vec{b}}{\vec{a}\T\vec{a}}\,\vec{a}$. The quantity
$\vec{a}\T\vec{b}$ is a number, so we are free to slide it around. Watch:
\[
  \vec{p} = \vec{a}\,\frac{\vec{a}\T\vec{b}}{\vec{a}\T\vec{a}}
          = \frac{\vec{a}\,\vec{a}\T}{\vec{a}\T\vec{a}}\,\vec{b}
          = \mat{P}\,\vec{b},
  \qquad\text{where}\qquad
  \mat{P} = \frac{\vec{a}\,\vec{a}\T}{\vec{a}\T\vec{a}}.
\]
Look at what happened. The object $\vec{a}\,\vec{a}\T$ is a \emph{column times a row} -- an
$n\times1$ times a $1\times n$ -- so it is an $n\times n$ \emph{matrix}, not a number. (Contrast
$\vec{a}\T\vec{a}$, a row times a column, which is a single number.) Dividing that matrix by the
number $\vec{a}\T\vec{a}$ gives the \textbf{projection matrix} $\mat{P}$. It does not depend on
$\vec{b}$ at all. Build it once, and then projecting \emph{any} vector onto the line is a single
matrix multiply $\vec{p}=\mat{P}\vec{b}$.

\begin{trapbox}
Keep the two products straight. $\vec{a}\T\vec{a}$ (row times column) is a \emph{scalar}; for
$\vec{a}=\colvec{2\\1}$ it is $5$. But $\vec{a}\,\vec{a}\T$ (column times row) is a \emph{matrix}:
\[
  \vec{a}\,\vec{a}\T = \colvec{2\\1}\begin{bmatrix}2 & 1\end{bmatrix}
  = \begin{bmatrix} 4 & 2 \\ 2 & 1\end{bmatrix}.
\]
Swapping them is the most common projection mistake. One collapses to a number; the other expands
to a square matrix. Order matters with $\T$.
\end{trapbox}

\begin{example}\label{ch18:line-Pmatrix}
For $\vec{a}=\colvec{2\\1}$ we found $\vec{a}\,\vec{a}\T=\begin{bmatrix}4&2\\2&1\end{bmatrix}$ and
$\vec{a}\T\vec{a}=5$, so
\[
  \mat{P} = \frac{1}{5}\begin{bmatrix} 4 & 2 \\ 2 & 1\end{bmatrix}
          = \begin{bmatrix} 0.8 & 0.4 \\ 0.4 & 0.2 \end{bmatrix}.
\]
Apply it to $\vec{b}=\colvec{1\\4}$ from Example~\ref{ch18:line-diagonal}:
\[
  \mat{P}\vec{b} = \frac{1}{5}\begin{bmatrix}4&2\\2&1\end{bmatrix}\colvec{1\\4}
  = \frac{1}{5}\colvec{4+8\\2+4} = \frac{1}{5}\colvec{12\\6} = \colvec{12/5\\6/5}.
\]
Same $\vec{p}$ as before. The matrix has packaged the whole projection into one reusable object.
\end{example}

This little $\mat{P}$ already shows the two features that turn out to define every projection
matrix. It is \textbf{symmetric}: $\mat{P}=\mat{P}\T$ (look at the $0.4$s off the diagonal). And
projecting twice does nothing new -- once you are on the line, projecting again leaves you put -- so
$\mat{P}^2=\mat{P}$. The second deserves a check:
\[
  \mat{P}^2 = \frac{1}{25}\begin{bmatrix}4&2\\2&1\end{bmatrix}\begin{bmatrix}4&2\\2&1\end{bmatrix}
  = \frac{1}{25}\begin{bmatrix}16+4 & 8+2\\ 8+2 & 4+1\end{bmatrix}
  = \frac{1}{25}\begin{bmatrix}20&10\\10&5\end{bmatrix}
  = \frac{1}{5}\begin{bmatrix}4&2\\2&1\end{bmatrix} = \mat{P}.
\]

\begin{aestheticbox}
Geometrically, an orthogonal projection retains the component of a vector lying
within a subspace and eliminates the component perpendicular to it. The algebraic
property $\mat{P}^2 = \mat{P}$ reflects this stabilization: once a vector has
been projected into a subspace $S$, applying the projection a second time leaves
the vector unchanged. A matrix satisfying $\mat{P}^2 = \mat{P}$ is called
\emph{idempotent}. As we will see when studying eigenvalues, an idempotent
matrix can scale eigenvectors only by $1$ (on $S$) or by $0$ (on $S^\perp$).
\end{aestheticbox}

\section{Why the perpendicular is the closest point}

We have used the fact that the perpendicular dropped to a subspace gives the nearest
point. We now give the formal proof using the Pythagorean theorem.

\begin{theorem}[The projection is the closest point]\label{ch18:closest-thm}
Let $\vec{p}$ be the projection of $\vec{b}$ onto a subspace $S$, so that the error
$\vec{e}=\vec{b}-\vec{p}$ is orthogonal to $S$. Then for \emph{every} other point $\vec{y}\in S$,
\[
  \norm{\vec{b}-\vec{y}} \;\ge\; \norm{\vec{b}-\vec{p}},
\]
with equality only when $\vec{y}=\vec{p}$. The projection is strictly the closest point of $S$ to
$\vec{b}$.
\end{theorem}

\begin{proof}
Pick any $\vec{y}\in S$ and split the trip from $\vec{b}$ to $\vec{y}$ into two legs through
$\vec{p}$:
\[
  \vec{b}-\vec{y} = (\vec{b}-\vec{p}) + (\vec{p}-\vec{y}).
\]
The first leg $\vec{b}-\vec{p}=\vec{e}$ is orthogonal to $S$. The second leg $\vec{p}-\vec{y}$
lives \emph{inside} $S$ (both $\vec{p}$ and $\vec{y}$ are in the subspace, and subspaces are closed
under subtraction). Orthogonal-to-$S$ and inside-$S$ means these two legs are perpendicular to each
other, so Pythagoras applies:
\[
  \norm{\vec{b}-\vec{y}}^2 = \norm{\vec{b}-\vec{p}}^2 + \norm{\vec{p}-\vec{y}}^2.
\]
The right side is $\norm{\vec{b}-\vec{p}}^2$ plus something that is never negative, so it is at
least $\norm{\vec{b}-\vec{p}}^2$. It equals it only when $\norm{\vec{p}-\vec{y}}^2=0$, i.e.\ when
$\vec{y}=\vec{p}$.
\end{proof}

\begin{figure}[H]
\centering
\begin{tikzpicture}[scale=1.3,>=Stealth]
  \draw[spacegray,thick] (-0.6,0)--(4.2,0);
  \node[spacegray,above] at (3.9,0.05){subspace $S$};
  % p on the line
  \fill (1.5,0) circle (1.4pt); \node[below] at (1.5,-0.08){$\vec{p}$};
  % y another point on the line
  \fill (3.2,0) circle (1.4pt); \node[below] at (3.2,-0.08){$\vec{y}$};
  % b above p
  \fill[vecred] (1.5,1.6) circle (1.4pt); \node[vecred,above] at (1.5,1.65){$\vec{b}$};
  % error leg (perpendicular)
  \draw[quotientorange,very thick] (1.5,1.6)--(1.5,0);
  \node[quotientorange,right] at (1.5,0.8){$\vec{e}$};
  % hypotenuse b to y
  \draw[vecblue,very thick] (1.5,1.6)--(3.2,0);
  \node[vecblue,above right] at (2.35,0.85){$\vec{b}-\vec{y}$};
  % in-subspace leg
  \draw[targetgreen,very thick] (1.5,0)--(3.2,0);
  \node[targetgreen,below] at (2.35,-0.12){$\vec{p}-\vec{y}$};
  % right angle
  \draw (1.5,0)++(0,0.16)--++(0.16,0)--++(0,-0.16);
\end{tikzpicture}
\caption{Why the perpendicular minimizes distance. Any other point $\vec{y}$ in $S$ sits at the far end of a
right triangle whose legs are the perpendicular error $\vec{e}$ and an in-subspace step. The
hypotenuse $\vec{b}-\vec{y}$ is strictly longer than the perpendicular leg $\vec{e}$ whenever $\vec{y}\neq\vec{p}$.}
\label{ch18:pythagoras}
\end{figure}

\begin{focusbox}
The entire theory of projection rests on a single geometric principle:
\textbf{the error vector is orthogonal to the subspace}. By the Pythagorean
theorem, the perpendicular distance from a point to a subspace is strictly
shorter than the distance to any other point in the subspace.
\end{focusbox}

\section{Projecting onto a subspace}

Now we leave the line behind. Suppose $S$ is not a single direction but a whole subspace -- a plane
in $\R^3$, or an $r$-dimensional flat slab inside $\R^n$. Give $S$ a set of basis vectors and stack
them as the columns of a matrix:
\[
  \mat{A} = \begin{bmatrix} \vec{a}_1 & \vec{a}_2 & \cdots & \vec{a}_r \end{bmatrix},
  \qquad
  S = \col(\mat{A}).
\]
Every point of $S$ is a combination of the columns, i.e.\ has the form $\mat{A}\vec{x}$ for some
coefficient vector $\vec{x}\in\R^r$. So projecting $\vec{b}$ onto $S$ means: find the coefficient
vector $\hat{\vec{x}}$ that makes $\vec{p}=\mat{A}\hat{\vec{x}}$ as close to $\vec{b}$ as possible.

\begin{intuitionbox}
For an $r$-dimensional subspace, we seek the coefficient vector
$\hat{\vec{x}} = (\hat{x}_1,\dots,\hat{x}_r)$ determining the optimal linear
combination $\vec{p} = \mat{A}\hat{\vec{x}} \in S$. As with a line, the condition
determining $\hat{\vec{x}}$ is that the error vector
$\vec{e} = \vec{b} - \mat{A}\hat{\vec{x}}$ must be orthogonal to every direction in $S$.
\end{intuitionbox}

\subsection{The normal equations}

The error is $\vec{e}=\vec{b}-\mat{A}\hat{\vec{x}}$. ``Perpendicular to $S$'' now means
perpendicular to \emph{every} column of $\mat{A}$, because the columns span $S$. (Why does checking
the basis columns suffice? Any vector in $S$ has the form $\mat{A}\vec{y}$; if $\mat{A}\T\vec{e}=\vec{0}$,
then the dot product with any vector in $S$ is $(\mat{A}\vec{y})\T\vec{e} = \vec{y}\T(\mat{A}\T\vec{e})
= \vec{y}\T\vec{0} = 0$, so $\vec{e}$ is automatically perpendicular to the entire subspace.)
So we need $\vec{a}_1\T\vec{e}=0$, $\vec{a}_2\T\vec{e}=0$, all the way to $\vec{a}_r\T\vec{e}=0$. Stacking
these $r$ equations is exactly what $\mat{A}\T\vec{e}=\vec{0}$ says:
\[
  \mat{A}\T\vec{e}
  = \begin{bmatrix} \vec{a}_1\T \\ \vdots \\ \vec{a}_r\T \end{bmatrix}\vec{e}
  = \begin{bmatrix} \vec{a}_1\T\vec{e} \\ \vdots \\ \vec{a}_r\T\vec{e} \end{bmatrix}
  = \vec{0}.
\]
Substitute $\vec{e}=\vec{b}-\mat{A}\hat{\vec{x}}$ and expand:
\[
  \mat{A}\T(\vec{b}-\mat{A}\hat{\vec{x}}) = \vec{0}
  \qquad\Longrightarrow\qquad
  \boxed{\;\mat{A}\T\mat{A}\,\hat{\vec{x}} = \mat{A}\T\vec{b}.\;}
\]
These are the famous \textbf{normal equations}. They are ``normal'' in the old sense of the word --
\emph{normal} means \emph{perpendicular}, and the whole system was born from the perpendicularity
condition $\mat{A}\T\vec{e}=\vec{0}$. Compare line and subspace side by side and feel how little
changed:
\[
  \text{line:}\quad (\vec{a}\T\vec{a})\,\hat{x} = \vec{a}\T\vec{b},
  \qquad\qquad
  \text{subspace:}\quad (\mat{A}\T\mat{A})\,\hat{\vec{x}} = \mat{A}\T\vec{b}.
\]
The scalar $\vec{a}\T\vec{a}$ grew into the $r\times r$ matrix $\mat{A}\T\mat{A}$; the number
$\vec{a}\T\vec{b}$ grew into the vector $\mat{A}\T\vec{b}$. Same skeleton, more dimensions.

\begin{formalbox}
\textbf{The projection matrix.} If the columns of $\mat{A}$ are independent, then $\mat{A}\T\mat{A}$
is invertible, and we may solve the normal equations:
\[
  \hat{\vec{x}} = (\mat{A}\T\mat{A})\inv\mat{A}\T\vec{b},
  \qquad
  \vec{p} = \mat{A}\hat{\vec{x}} = \mat{A}(\mat{A}\T\mat{A})\inv\mat{A}\T\,\vec{b}.
\]
The matrix that projects any $\vec{b}$ onto $\col(\mat{A})$ is therefore
\[
  \boxed{\;\mat{P} = \mat{A}(\mat{A}\T\mat{A})\inv\mat{A}\T.\;}
\]
\end{formalbox}

Notice how this collapses to the line formula when $\mat{A}$ is a single column $\vec{a}$. Then
$\mat{A}\T\mat{A}=\vec{a}\T\vec{a}$ is the scalar we divide by, $(\mat{A}\T\mat{A})\inv$ is just
$1/(\vec{a}\T\vec{a})$, and $\mat{P}=\vec{a}(\vec{a}\T\vec{a})\inv\vec{a}\T
= \dfrac{\vec{a}\,\vec{a}\T}{\vec{a}\T\vec{a}}$. The grand formula and the baby formula are the same
formula.

\begin{trapbox}
You might be tempted to ``simplify''
$\mat{A}(\mat{A}\T\mat{A})\inv\mat{A}\T$ by distributing the inverse:
$\mat{A}\mat{A}\inv(\mat{A}\T)\inv\mat{A}\T = \mat{I}$. \textbf{Do not.} That algebra is illegal
unless $\mat{A}$ is square and invertible -- and if $\mat{A}$ were square and invertible, its column
space would be \emph{all} of $\R^n$, every $\vec{b}$ would already be in $S$, and projection would
indeed be the identity. The interesting case is the rectangular one ($m>r$): $\mat{A}$ has more
rows than columns, $\mat{A}\inv$ does not exist, and $\mat{P}$ is genuinely \emph{not} the
identity. Treat $(\mat{A}\T\mat{A})\inv$ as one indivisible block.
\end{trapbox}

\begin{recoverybox}
\textbf{Mental check: when $\mat{A}$ is square and invertible.} What happens to the formula
$\mat{P} = \mat{A}(\mat{A}\T\mat{A})\inv\mat{A}\T$ if $\mat{A}$ is a square, invertible $n\times n$ matrix?
\emph{Pause and predict before checking the algebra:} If the columns of $\mat{A}$ form a basis of all of
$\R^n$, then $\col(\mat{A})$ is the entire ambient space. Every vector $\vec{b}$ already lives in the
subspace, so the closest point to $\vec{b}$ is $\vec{b}$ itself, and the projection must be the
identity matrix $\mat{I}$. And indeed, only in this square invertible case can we legally distribute the
inverse: $\mat{A}(\mat{A}\T\mat{A})\inv\mat{A}\T = \mat{A}\mat{A}\inv(\mat{A}\T)\inv\mat{A}\T = \mat{I}\cdot\mat{I} = \mat{I}$.
\end{recoverybox}

\subsection{Why $\mat{A}\T\mat{A}$ is invertible when the columns are independent}

We leaned on ``$\mat{A}\T\mat{A}$ is invertible'' -- let us prove it now. It is a small gem
of an argument and it tells you exactly when projection works.

\begin{lemma}\label{ch18:ata-invertible}
If the columns of $\mat{A}$ are linearly independent, then $\mat{A}\T\mat{A}$ is invertible.
\end{lemma}

\begin{proof}
$\mat{A}\T\mat{A}$ is a square $r\times r$ matrix, so it suffices to show its nullspace is trivial:
the only $\vec{x}$ with $\mat{A}\T\mat{A}\vec{x}=\vec{0}$ is $\vec{x}=\vec{0}$. Suppose
$\mat{A}\T\mat{A}\vec{x}=\vec{0}$. Dot both sides with $\vec{x}$:
\[
  \vec{x}\T\mat{A}\T\mat{A}\vec{x} = 0
  \qquad\Longrightarrow\qquad
  (\mat{A}\vec{x})\T(\mat{A}\vec{x}) = 0
  \qquad\Longrightarrow\qquad
  \norm{\mat{A}\vec{x}}^2 = 0.
\]
A vector with zero length is the zero vector, so $\mat{A}\vec{x}=\vec{0}$. But the columns of
$\mat{A}$ are independent, which means the \emph{only} combination of them giving zero is the
trivial one: $\vec{x}=\vec{0}$. Hence $\mat{A}\T\mat{A}$ has trivial nullspace and is invertible.
\end{proof}

The middle step is a trick worth keeping: $\vec{x}\T\mat{A}\T\mat{A}\vec{x}=\norm{\mat{A}\vec{x}}^2$
turns a matrix expression into a squared length, which can only be zero one way. You will see this
move again in the chapter on positive-definite matrices, where it is the central theme.

\section{The two properties that define a projection}

Whatever $\mat{A}$ is, the matrix $\mat{P}=\mat{A}(\mat{A}\T\mat{A})\inv\mat{A}\T$ has two
properties, and they are so characteristic that they \emph{define} what it means to be an
orthogonal projection.

\begin{theorem}[Symmetry and idempotence]\label{ch18:P-props}
The projection matrix $\mat{P}=\mat{A}(\mat{A}\T\mat{A})\inv\mat{A}\T$ satisfies
\[
  \text{(1)}\quad \mat{P}\T = \mat{P}
  \qquad\text{and}\qquad
  \text{(2)}\quad \mat{P}^2 = \mat{P}.
\]
\end{theorem}

\begin{proof}
\emph{Symmetry.} Transpose, remembering that $(\mat{XYZ})\T=\mat{Z}\T\mat{Y}\T\mat{X}\T$ and that
$\mat{A}\T\mat{A}$ is itself symmetric (so its inverse is too):
\[
  \mat{P}\T = \big(\mat{A}(\mat{A}\T\mat{A})\inv\mat{A}\T\big)\T
  = (\mat{A}\T)\T\,\big((\mat{A}\T\mat{A})\inv\big)\T\,\mat{A}\T
  = \mat{A}\,(\mat{A}\T\mat{A})\inv\,\mat{A}\T = \mat{P}.
\]
\emph{Idempotence.} Multiply $\mat{P}$ by itself and watch the middle collapse. Let
$\mat{M}=(\mat{A}\T\mat{A})\inv$ for brevity:
\[
  \mat{P}^2 = \mat{A}\mat{M}\mat{A}\T\,\mat{A}\mat{M}\mat{A}\T
  = \mat{A}\mat{M}\,(\mat{A}\T\mat{A})\,\mat{M}\,\mat{A}\T.
\]
But $\mat{M}(\mat{A}\T\mat{A}) = (\mat{A}\T\mat{A})\inv(\mat{A}\T\mat{A}) = \mat{I}$, so the middle
three factors $\mat{M}(\mat{A}\T\mat{A})\mat{M}$ collapse to $\mat{M}$, leaving
$\mat{P}^2 = \mat{A}\mat{M}\mat{A}\T = \mat{P}$.
\end{proof}

\begin{intuitionbox}
Both properties are obvious \emph{geometrically}, which is the best kind of proof to carry in your
head. Idempotence $\mat{P}^2=\mat{P}$: once $\vec{b}$'s shadow has landed in $S$, it is already in
$S$, so projecting it again does nothing. Symmetry $\mat{P}\T=\mat{P}$ is subtler but says the
projection ``treats both directions fairly'': the part of $\vec{b}$ that points into $\vec{c}$'s
shadow equals the part of $\vec{c}$ that points into $\vec{b}$'s shadow, i.e.\
$\vec{c}\T\mat{P}\vec{b} = (\mat{P}\vec{c})\T\vec{b}$. Drop symmetry and you get a \emph{slanted}
(oblique) projection that still satisfies $\mat{P}^2=\mat{P}$ but no longer drops a true
perpendicular -- and no longer finds the closest point (Figure~\ref{ch18:orthog-vs-oblique}).
\end{intuitionbox}

\begin{figure}[H]
\centering
\begin{tikzpicture}[scale=1.1,>=Stealth]
  % Left panel: Orthogonal Projection
  \begin{scope}[shift={(0,0)}]
    \draw[spacegray,thick] (-0.5,0) -- (3.5,0);
    \node[spacegray,below] at (3.2,-0.1) {$S$};
    \node[align=center] at (1.5,2.5) {\small\bfseries Orthogonal Projection\\[2pt]\scriptsize\color{spacegray}$\mat{P}^2 = \mat{P}$ and $\mat{P}\T = \mat{P}$};
    
    % point b
    \fill[vecred] (1.5,1.7) circle (1.5pt);
    \node[vecred,above] at (1.5,1.75) {$\vec{b}$};
    
    % projection p
    \fill[targetgreen] (1.5,0) circle (1.5pt);
    \node[targetgreen,below] at (1.5,-0.1) {$\vec{p}$};
    
    % error vector (straight down)
    \draw[->,quotientorange,very thick] (1.5,1.7) -- (1.5,0.05);
    \node[quotientorange,right,font=\scriptsize] at (1.5,0.85) {$\vec{e} \perp S$};
    
    % right angle mark
    \draw (1.5,0) ++(0,0.16) -- ++(0.16,0) -- ++(0,-0.16);
    \node[font=\scriptsize,targetgreen!70!black] at (1.5,-0.75) {closest point on $S$};
  \end{scope}

  % Right panel: Oblique Projection
  \begin{scope}[shift={(5.2,0)}]
    \draw[spacegray,thick] (-0.5,0) -- (3.5,0);
    \node[spacegray,below] at (3.2,-0.1) {$S$};
    \node[align=center] at (1.5,2.5) {\small\bfseries Oblique Projection\\[2pt]\scriptsize\color{spacegray}$\mat{P}^2 = \mat{P}$ but not symmetric};
    
    % point b
    \fill[vecred] (1.2,1.7) circle (1.5pt);
    \node[vecred,above] at (1.2,1.75) {$\vec{b}$};
    
    % projection p_oblique
    \fill[targetgreen] (2.4,0) circle (1.5pt);
    \node[targetgreen,below] at (2.4,-0.1) {$\vec{p}_{\text{oblique}}$};
    
    % error vector (slanted)
    \draw[->,quotientorange,very thick] (1.2,1.7) -- (2.36,0.05);
    \node[quotientorange,above right,font=\scriptsize] at (1.8,0.85) {slanted error};
    
    % perpendicular foot for comparison (ghost)
    \draw[spacegray,dashed] (1.2,1.7) -- (1.2,0);
    \fill[spacegray] (1.2,0) circle (1.0pt);
    \node[spacegray,above left,font=\tiny] at (1.18,0.05) {nearest};
    
    \node[font=\scriptsize,quotientorange!80!black] at (1.5,-0.75) {idempotent, but not closest};
  \end{scope}
\end{tikzpicture}
\caption{Why symmetry matters. \textbf{Left:} when $\mat{P}$ is both idempotent and symmetric, the error $\vec{e}$ drops straight down at a right angle, finding the unique nearest point. \textbf{Right:} when $\mat{P}^2 = \mat{P}$ but $\mat{P}\T \ne \mat{P}$, the shadow slides along a tilted ray. It is still a projection (once on $S$, projecting again does nothing), but it fails to minimize the distance to $\vec{b}$.}
\label{ch18:orthog-vs-oblique}
\end{figure}

\begin{recoverybox}
\textbf{Mental check: projecting the zero vector.} Where does any projection matrix send $\vec{0}$?
And if a matrix satisfies $\mat{P}\vec{v} = \vec{v}$ for \emph{every} vector $\vec{v}$ in $\R^n$, what is $\mat{P}$?
\emph{Pause and think:} $\mat{P}\vec{0} = \vec{0}$, because every linear subspace passes through the
origin. If $\mat{P}\vec{v}=\vec{v}$ everywhere, the subspace is the entire ambient space $\R^n$, so the error is
always zero and $\mat{P} = \mat{I}$.
\end{recoverybox}

\begin{historybox}
The geometric idea that the shortest path from a point to a plane is the perpendicular dropped onto
it was known to ancient Greek geometers. But converting that geometric intuition into an explicit
algebraic operator $\mat{P}$ that acts on arbitrary multidimensional vectors was the achievement of
nineteenth- and twentieth-century analysis.

The concept matured through physical variational principles, notably Dirichlet's principle in potential
theory, where nature minimizes field energy by finding the projection of a state onto admissible
boundary configurations. In the early 1900s, Maurice Fréchet and David Hilbert formalized projection
theorems in general inner-product spaces. They proved that in any complete linear space, every vector
decomposes uniquely into an orthogonal projection and a perpendicular error vector that achieves the
absolute minimum Euclidean distance.
\end{historybox}

\begin{remark}
The complementary matrix $\mat{I}-\mat{P}$ is \emph{also} a projection -- it projects onto the
orthogonal complement of $S$, the space where the errors live. Check:
$(\mat{I}-\mat{P})^2 = \mat{I}-2\mat{P}+\mat{P}^2 = \mat{I}-2\mat{P}+\mat{P} = \mat{I}-\mat{P}$,
and it is symmetric because $\mat{P}$ is. So $\mat{P}$ keeps the part of $\vec{b}$ inside $S$, and
$\mat{I}-\mat{P}$ keeps the part orthogonal to $S$; together they split every vector cleanly,
$\vec{b} = \mat{P}\vec{b} + (\mat{I}-\mat{P})\vec{b} = \vec{p} + \vec{e}$.
\end{remark}

\begin{figure}[H]
\centering
\begin{tikzpicture}[scale=1.3,>=Stealth]
  % subspace S as a horizontal line, S-perp as vertical
  \draw[targetgreen,thick] (-2.6,0)--(2.6,0);
  \node[targetgreen,right] at (2.6,0){$S$ \small($\mat{P}\vec{v} = \vec{v}$, kept)};
  \draw[structurepurple,thick] (0,-0.4)--(0,2.6);
  \node[structurepurple,above] at (0,2.6){$S^\perp$ \small($\mat{P}\vec{v} = \vec{0}$, crushed)};
  % b
  \draw[->,vecred,very thick] (0,0)--(2.0,1.8);
  \node[vecred,above right] at (2.0,1.8){$\vec{b}$};
  % p = P b on S
  \draw[->,targetgreen!60!black,very thick] (0,0)--(2.0,0);
  \node[targetgreen!60!black,below] at (1.0,0){$\vec{p}=\mat{P}\vec{b}$};
  % e = (I-P) b on S-perp
  \draw[->,quotientorange,very thick] (0,0)--(0,1.8);
  \node[quotientorange,left] at (0,1.0){$\vec{e}=(\mat{I}-\mat{P})\vec{b}$};
  % dashed parallelogram closure
  \draw[spacegray,dashed] (2.0,0)--(2.0,1.8);
  \draw[spacegray,dashed] (0,1.8)--(2.0,1.8);
  % right angle at origin
  \draw (0,0)++(0.16,0)--++(0,0.16)--++(-0.16,0);
\end{tikzpicture}
\caption{One projection splits space in two. The kept part $\vec{p}=\mat{P}\vec{b}$ lives in $S$
(where $\mat{P}$ acts as the identity: $\mat{P}\vec{p}=\vec{p}$); the forgotten part
$\vec{e}=(\mat{I}-\mat{P})\vec{b}$ lives in the orthogonal complement $S^\perp$ (where $\mat{P}$
sends everything to zero: $\mat{P}\vec{e}=\vec{0}$). Their right-angled sum rebuilds $\vec{b}$.
(In Chapter~23, these invariant directions will be recognized as eigenspaces for scale factors $1$ and $0$.)}
\label{ch18:decomp-fig}
\end{figure}

\section{A fully worked plane in three dimensions}

To see the projection formula in concrete computation, let us project a point onto a plane
in $\R^3$ from start to finish.

\begin{example}\label{ch18:plane-worked}
Project $\vec{b}=\colvec{6\\0\\0}$ onto the plane $S$ spanned by
$\vec{a}_1=\colvec{1\\1\\0}$ and $\vec{a}_2=\colvec{0\\1\\1}$. Assemble
\[
  \mat{A} = \begin{bmatrix} 1 & 0 \\ 1 & 1 \\ 0 & 1 \end{bmatrix}.
\]
\textbf{Step 1: form $\mat{A}\T\mat{A}$.}
\[
  \mat{A}\T\mat{A}
  = \begin{bmatrix} 1 & 1 & 0 \\ 0 & 1 & 1 \end{bmatrix}
    \begin{bmatrix} 1 & 0 \\ 1 & 1 \\ 0 & 1 \end{bmatrix}
  = \begin{bmatrix} 2 & 1 \\ 1 & 2 \end{bmatrix}.
\]
\textbf{Step 2: form $\mat{A}\T\vec{b}$.}
\[
  \mat{A}\T\vec{b}
  = \begin{bmatrix} 1 & 1 & 0 \\ 0 & 1 & 1 \end{bmatrix}\colvec{6\\0\\0}
  = \colvec{6\\0}.
\]
\textbf{Step 3: solve the normal equations $\mat{A}\T\mat{A}\,\hat{\vec{x}}=\mat{A}\T\vec{b}$.}
The inverse of $\begin{bmatrix}2&1\\1&2\end{bmatrix}$ is
$\frac{1}{3}\begin{bmatrix}2&-1\\-1&2\end{bmatrix}$ (determinant $4-1=3$), so
\[
  \hat{\vec{x}} = \frac{1}{3}\begin{bmatrix}2&-1\\-1&2\end{bmatrix}\colvec{6\\0}
  = \frac{1}{3}\colvec{12\\-6} = \colvec{4\\-2}.
\]
\textbf{Step 4: form the projection $\vec{p}=\mat{A}\hat{\vec{x}}$.}
\[
  \vec{p} = \begin{bmatrix}1&0\\1&1\\0&1\end{bmatrix}\colvec{4\\-2}
  = \colvec{4\\\,4-2\,\\-2} = \colvec{4\\2\\-2}.
\]
\textbf{Step 5: the error and the all-important check.}
\[
  \vec{e} = \vec{b}-\vec{p} = \colvec{6\\0\\0}-\colvec{4\\2\\-2} = \colvec{2\\-2\\2}.
\]
Is $\vec{e}$ perpendicular to the plane? Dot it with each basis vector:
\[
  \vec{a}_1\T\vec{e} = 1\cdot2 + 1\cdot(-2) + 0\cdot2 = 0,
  \qquad
  \vec{a}_2\T\vec{e} = 0\cdot2 + 1\cdot(-2) + 1\cdot2 = 0.
\]
Both zero. The error is genuinely orthogonal to the plane, so $\vec{p}=\colvec{4\\2\\-2}$ is the
closest point of $S$ to $\vec{b}$. Notice $\vec{e}=\colvec{2\\-2\\2}=2\colvec{1\\-1\\1}$ points along
the plane's \emph{normal} direction $\colvec{1\\-1\\1}$ -- exactly the direction perpendicular to
both $\vec{a}_1$ and $\vec{a}_2$, as it must.
\end{example}

\begin{figure}[H]
\centering
\begin{tikzpicture}[scale=1.1,>=Stealth]
  % Plane S spanned by a1 and a2
  \coordinate (O) at (0,0);
  \fill[targetgreen!12] (-2.5,-1.1) -- (2.7,-0.7) -- (3.5,1.5) -- (-1.7,1.1) -- cycle;
  \draw[targetgreen!70!black,thick] (-2.5,-1.1) -- (2.7,-0.7) -- (3.5,1.5) -- (-1.7,1.1) -- cycle;
  \node[draw=targetgreen!80, fill=targetgreen!15, rounded corners=2pt, font=\scriptsize\bfseries, text=targetgreen!95!black]
    at (2.7,0.8) {Subspace $S = \col(\mat{A})$};

  % Origin
  \fill[black] (O) circle (1.8pt) node[below left=2pt, font=\scriptsize] {$\vec{0}$};

  % Basis vectors in the plane
  \coordinate (A1) at (1.8,-0.3);
  \coordinate (A2) at (-0.8,0.7);
  \draw[->,very thick,vecblue] (O) -- (A1) node[below right=1pt, font=\scriptsize] {$\vec{a}_1$};
  \draw[->,very thick,vecblue] (O) -- (A2) node[above left=1pt, font=\scriptsize] {$\vec{a}_2$};

  % Projection p in S
  \coordinate (P) at (1.2,0.4);
  \draw[->,very thick,targetgreen!85!black] (O) -- (P);
  \node[targetgreen!85!black, font=\scriptsize, below right=1pt] at (P) {$\vec{p} = \mat{A}\hat{\vec{x}}$};
  \fill[targetgreen!85!black] (P) circle (1.8pt);

  % Vector b off the plane
  \coordinate (B) at (1.2,2.6);
  \draw[->,very thick,vecred] (O) -- (B) node[above right=1pt, font=\scriptsize] {$\vec{b}$};
  \fill[vecred] (B) circle (1.8pt);

  % Perpendicular error vector e = b - p
  \draw[->,line width=1.6pt,quotientorange!95!black] (P) -- (B)
    node[midway,right=3pt,font=\small,fill=white,inner sep=1.5pt] {$\vec{e} = \vec{b} - \vec{p}$};

  % Right angle indicator at P
  \draw[quotientorange!80!black,thick] (1.2,0.65) -- (0.95,0.65) -- (0.95,0.4);

  % Normal line extending through error
  \draw[dashed,spacegray,thick] (B) -- (1.2,3.3) node[above, font=\scriptsize] {$\vec{n} \in S^\perp$};

  % Orthogonality annotation below plane
  \node[font=\scriptsize, spacegray, align=center] at (0.4,-1.55) {Error $\vec{e} \perp S \implies \mat{A}\T\vec{e} = \vec{0} \implies \mat{A}\T\mat{A}\hat{\vec{x}} = \mat{A}\T\vec{b}$};
\end{tikzpicture}
\caption{Orthogonal projection of $\vec{b}$ onto a 2D subspace $S = \col(\mat{A})$.
The projection $\vec{p} = \mat{A}\hat{\vec{x}}$ is the unique closest point in the plane.
The residual error vector $\vec{e} = \vec{b} - \vec{p}$ drops straight down, forming a right angle
with the basis directions $\vec{a}_1, \vec{a}_2$ and therefore with every vector in $S$.}
\label{ch18:plane-fig}
\end{figure}

\begin{labbox}
Build the full $3\times3$ projection matrix for that plane and feel it work. With
$\mat{A}\T\mat{A}=\begin{bmatrix}2&1\\1&2\end{bmatrix}$ and inverse
$\frac13\begin{bmatrix}2&-1\\-1&2\end{bmatrix}$, multiply out
$\mat{P}=\mat{A}(\mat{A}\T\mat{A})\inv\mat{A}\T$. You should land on
\[
  \mat{P} = \frac{1}{3}\begin{bmatrix} 2 & 1 & -1 \\ 1 & 2 & 1 \\ -1 & 1 & 2 \end{bmatrix}.
\]
Check three things by hand: (a) $\mat{P}$ is symmetric; (b) $\mat{P}\vec{b}=\colvec{4\\2\\-2}$ for
$\vec{b}=\colvec{6\\0\\0}$, matching Step~4; (c) the \emph{trace} (the sum of the diagonal entries,
$\tr(\mat{P}) = P_{11}+P_{22}+P_{33}$) is $\frac13(2+2+2)=2$, which directly equals the dimension
of the plane. Summing the diagonal of any projection matrix always counts the dimension of the subspace
it projects onto.
\end{labbox}

\section{Two special cases worth checking}

Two extreme cases test whether you really understand $\mat{P}=\mat{A}(\mat{A}\T\mat{A})\inv\mat{A}\T$.

\textbf{When $\vec{b}$ is already in the space.} If $\vec{b}\in\col(\mat{A})$ then its closest point
is itself, so we expect $\mat{P}\vec{b}=\vec{b}$. Indeed, $\vec{b}=\mat{A}\vec{c}$ for some
$\vec{c}$, and
\[
  \mat{P}\vec{b} = \mat{A}(\mat{A}\T\mat{A})\inv\mat{A}\T\mat{A}\vec{c}
  = \mat{A}(\mat{A}\T\mat{A})\inv(\mat{A}\T\mat{A})\vec{c}
  = \mat{A}\vec{c} = \vec{b}.
\]
Projection leaves the space it projects onto fixed.

\textbf{When $\vec{b}$ is orthogonal to the space.} If $\vec{b}\perp\col(\mat{A})$ then
$\mat{A}\T\vec{b}=\vec{0}$ (every column dots to zero with $\vec{b}$), so $\hat{\vec{x}}=\vec{0}$
and $\vec{p}=\vec{0}$. The shadow collapses to the origin. All of $\vec{b}$ becomes error.

\begin{aestheticbox}
Compare how $\mat{P}$ acts on these two families of vectors.
Vectors \emph{in} the subspace $S$ are untouched: $\mat{P}\vec{v} = 1\cdot\vec{v}$ (they are kept).
Vectors \emph{perpendicular} to the subspace are mapped to zero: $\mat{P}\vec{v} = 0\cdot\vec{v}$.

This is our first encounter with one of the most central themes of linear algebra: finding the special directions that a matrix merely stretches by a numerical scale factor rather than rotating. In Chapter~23, we will formalize these directions as \emph{eigenvectors}, and their scaling factors as \emph{eigenvalues}. A projection matrix has only two possible scaling factors: $1$ (for directions inside $S$) and $0$ (for directions perpendicular to $S$). The algebraic condition $\mat{P}^2=\mat{P}$ demands this: if scaling a vector twice has the exact same effect as scaling it once ($\lambda^2=\lambda$), the scale factor $\lambda$ can only ever be $1$ or $0$. The number of independent directions scaled by $1$ is the dimension of $S$---which also equals the rank of $\mat{P}$, and the sum of its diagonal entries (the trace). The spectral structure of an orthogonal projection is entirely determined by these two eigenvalues: $1$ and $0$.
\end{aestheticbox}

\section{Where this is heading: the inconsistent system}

We opened with the promise that projection rescues the unsolvable equation $\mat{A}\vec{x}=\vec{b}$.
Let us close the loop. When $\vec{b}$ is not in the column space, no $\vec{x}$ makes
$\mat{A}\vec{x}=\vec{b}$. The sensible move is to solve instead the \emph{nearest} solvable problem:
replace $\vec{b}$ by its projection $\vec{p}$ onto $\col(\mat{A})$, and solve $\mat{A}\hat{\vec{x}}
=\vec{p}$. By construction $\vec{p}$ \emph{is} in the column space, so this has a solution -- and it
is the very $\hat{\vec{x}}$ from the normal equations. We have, almost without trying, derived the
least-squares solution.

\begin{focusbox}
\textbf{The bridge to least squares.} To ``solve'' an inconsistent $\mat{A}\vec{x}=\vec{b}$, project
$\vec{b}$ onto $\col(\mat{A})$ and accept the projection $\vec{p}$ as the best we can hit. The
best-fit coefficients are
\[
  \hat{\vec{x}} = (\mat{A}\T\mat{A})\inv\mat{A}\T\vec{b},
\]
and they minimize the squared error $\norm{\mat{A}\vec{x}-\vec{b}}^2$ -- because, by
Theorem~\ref{ch18:closest-thm}, the projection is the closest point. The next chapter puts this
one idea to work on curve fitting and regression.
\end{focusbox}

\begin{applicationbox}
\textbf{Least squares is everywhere a model meets noisy data.} Suppose you measure a quantity that
you believe depends linearly on inputs -- house price versus square footage and number of bedrooms,
say. Each data point is a row of $\mat{A}$ (the features) and an entry of $\vec{b}$ (the observed
price). There are far more houses than features, so $\mat{A}\vec{x}=\vec{b}$ is wildly
overdetermined and has no exact solution: no single line passes through every messy point. The
linear-regression coefficients that every spreadsheet, statistics package, and machine-learning
library computes are \emph{exactly} $\hat{\vec{x}}=(\mat{A}\T\mat{A})\inv\mat{A}\T\vec{b}$. The
``predictions'' $\hat{\vec{b}}=\mat{P}\vec{b}$ are the projection of reality onto the span of your
features, and the residuals $\vec{e}=\vec{b}-\hat{\vec{b}}$ are the part of reality your model
cannot explain -- guaranteed orthogonal to every feature, which is why a well-fit linear model's
residuals show no leftover linear trend. The matrix $\mat{P}$ even has a name in statistics: the
\emph{hat matrix}, because it puts the hat on $\vec{b}$ to make $\hat{\vec{b}}$.
\end{applicationbox}

\begin{connectionbox}
\textbf{Projection in function spaces: Fourier series.} The same perpendicular-drop works when the
``vectors'' are functions. Take the inner product $\inner{f}{g}=\int_0^{2\pi} f(t)g(t)\,dt$. The
functions $1,\cos t,\sin t,\cos 2t,\sin 2t,\dots$ are mutually orthogonal under this inner product
-- just like an orthogonal basis of $\R^n$, but infinite. A Fourier series approximation of a
signal $f$ is precisely the \emph{projection} of $f$ onto the finite-dimensional subspace spanned
by the first few of these waves. Each Fourier coefficient is a one-dimensional projection,
$\hat{x}_k = \inner{f}{\phi_k}/\inner{\phi_k}{\phi_k}$ -- the very formula from
Section~\ref{ch18:project-line}, now with integrals for dot products. ``Find the closest combination
of sines and cosines'' \emph{is} least squares in a space of functions, and the leftover (the part
$f$ minus its truncated series) is orthogonal to every wave you kept. This is the doorway from
linear algebra into Fourier analysis and the whole field of signal processing.
\end{connectionbox}

\begin{strangbox}
Strang builds projections in essentially this order -- line first, then subspace, then the matrix
$\mat{P}=\mat{A}(\mat{A}\T\mat{A})\inv\mat{A}\T$ -- and he hammers the same two properties,
$\mat{P}\T=\mat{P}$ and $\mat{P}^2=\mat{P}$. If you are reading this \emph{before} Strang, the move
to internalize is ``error orthogonal to the space $\Rightarrow$ closest point $\Rightarrow$ normal
equations.'' Everything in his projection and least-squares lectures hangs on that chain. If you are
reading this \emph{after} Strang, look ahead: $\mat{P}$ is your first nontrivial example of a matrix
whose action is completely described by its eigenvalues ($0$ and $1$) and eigenspaces ($S$ and
$S^\perp$). When you reach the spectral theorem and the SVD, you will recognize projection as the
simplest member of the family of symmetric matrices, and you will write any symmetric matrix as a
weighted sum of projections onto its eigenvectors.
\end{strangbox}

\section{What to carry forward}

\begin{focusbox}
\textbf{Summary of core concepts.}
\begin{itemize}
  \item To find the closest point of a subspace $S$ to $\vec{b}$, drop a perpendicular: demand
        that the error $\vec{e}=\vec{b}-\vec{p}$ be orthogonal to $S$.
  \item \textbf{Onto a line} $\spn(\vec{a})$:
        $\;\vec{p}=\dfrac{\vec{a}\T\vec{b}}{\vec{a}\T\vec{a}}\,\vec{a}$,
        \; $\mat{P}=\dfrac{\vec{a}\,\vec{a}\T}{\vec{a}\T\vec{a}}$.
  \item \textbf{Onto a subspace} $\col(\mat{A})$: solve the normal equations
        $\mat{A}\T\mat{A}\,\hat{\vec{x}}=\mat{A}\T\vec{b}$, then
        $\mat{P}=\mat{A}(\mat{A}\T\mat{A})\inv\mat{A}\T$.
  \item Every orthogonal projection matrix is \textbf{symmetric} ($\mat{P}\T=\mat{P}$) and \textbf{idempotent}
        ($\mat{P}^2=\mat{P}$). It acts as the identity on $S$ ($\mat{P}\vec{v}=\vec{v}$), collapses
        $S^\perp$ to the origin ($\mat{P}\vec{v}=\vec{0}$), and has $\rank(\mat{P}) = \dim S$.
  \item The closest-point guarantee (Pythagoras) is what turns projection into \textbf{least
        squares}, the engine of regression and data fitting.
\end{itemize}
\end{focusbox}

\begin{connectionbox}
Projection matrices replace impossible equations with optimal approximations. Whenever $\mat{A}\vec{x}=\vec{b}$ has no solution because $\vec{b} \notin \col(\mat{A})$, the projection matrix $\mat{P} = \mat{A}(\mat{A}\T\mat{A})\inv\mat{A}\T$ cleanly splits $\vec{b}$ into its reachable shadow $\vec{p} \in \col(\mat{A})$ and its orthogonal residual $\vec{e} \in \nullsp(\mat{A}\T)$.

Yet computing $\mat{P}$ directly by inverting $\mat{A}\T\mat{A}$ carries a severe hidden risk: forming the normal equations squares the matrix condition number $\kappa(\mat{A})$ (which measures sensitivity to numerical rounding and noise), yielding $\kappa(\mat{A}\T\mat{A}) = (\kappa(\mat{A}))^2$. If the columns of $\mat{A}$ are nearly parallel, $\mat{A}\T\mat{A}$ becomes dangerously ill-conditioned. In Chapter~19, we apply projections to data fitting and linear regression. Then, in Chapter~20, we resolve the numerical danger of $\mat{A}\T\mat{A}$ by developing the Gram--Schmidt process and QR factorization, computing orthogonal projections with pristine numerical stability.
\end{connectionbox}

\section{Exercises}

\begin{exercise}\label{ch18:ex-mental-checks}
\textbf{(Mental checks: three instant projections.)} Answer each of these without computing any matrix inverses:
\begin{enumerate}[label=(\alph*)]
  \item What is the $3\times3$ matrix $\mat{P}$ that projects every vector in $\R^3$ onto the $xy$-plane? (Think: which coordinates are kept, and which are dropped?)
  \item If $\vec{b}$ lies on the line spanned by $\vec{a}$, what is its projection $\mat{P}\vec{b}$? If $\vec{b}$ is perpendicular to $\vec{a}$, what is $\mat{P}\vec{b}$?
  \item If an orthogonal projection matrix $\mat{P}$ satisfies $\mat{P}^2=\mat{P}$ and $\mat{P}\T=\mat{P}$, what is $(\mat{I}-\mat{P})^2$? What is $\mat{P}(\mat{I}-\mat{P})$?
\end{enumerate}
\end{exercise}

\begin{exercise}\label{ch18:ex-line-basic}
Project $\vec{b}=\colvec{3\\4}$ onto the line through $\vec{a}=\colvec{1\\1}$. Find $\hat{x}$,
$\vec{p}$, and $\vec{e}$, and verify that $\vec{a}\T\vec{e}=0$.
\end{exercise}

\begin{exercise}\label{ch18:ex-Pmatrix}
For $\vec{a}=\colvec{3\\4}$, build the line-projection matrix
$\mat{P}=\dfrac{\vec{a}\,\vec{a}\T}{\vec{a}\T\vec{a}}$ explicitly as a $2\times2$ matrix. Confirm by
direct multiplication that $\mat{P}^2=\mat{P}$ and that $\mat{P}\T=\mat{P}$.
\end{exercise}

\begin{exercise}\label{ch18:ex-already-in}
Let $\mat{P}$ be the matrix from Exercise~\ref{ch18:ex-Pmatrix}. Compute $\mat{P}\colvec{3\\4}$ and
explain geometrically why the answer is what it is. Then compute $\mat{P}\colvec{4\\-3}$ and explain
that answer too. (Hint: how is $\colvec{4\\-3}$ related to $\colvec{3\\4}$?)
\end{exercise}

\begin{exercise}\label{ch18:ex-plane}
Project $\vec{b}=\colvec{0\\0\\3}$ onto the plane spanned by $\vec{a}_1=\colvec{1\\0\\1}$ and
$\vec{a}_2=\colvec{0\\1\\1}$. Form $\mat{A}\T\mat{A}$ and $\mat{A}\T\vec{b}$, solve for
$\hat{\vec{x}}$, find $\vec{p}$ and $\vec{e}$, and check that $\vec{e}$ is orthogonal to both basis
vectors. (You should find $\hat{\vec{x}}=\colvec{1\\1}$, $\vec{p}=\colvec{1\\1\\2}$, and
$\vec{e}=\colvec{-1\\-1\\1}$; note that $\vec{b}$ does \emph{not} lie in the plane, so the projection
genuinely moves it.)
\end{exercise}

\begin{exercise}\label{ch18:ex-trace}
Recall that the \emph{trace} of a square matrix is the sum of its diagonal entries, $\tr(\mat{M}) = \sum M_{ii}$,
and that the trace is cyclic: $\tr(\mat{X}\mat{Y})=\tr(\mat{Y}\mat{X})$ whenever both products are defined.
Show that for the plane in Exercise~\ref{ch18:ex-plane} the trace of the projection matrix
$\mat{P}=\mat{A}(\mat{A}\T\mat{A})\inv\mat{A}\T$ equals $2$, without computing all nine entries.
(Hint: let $\mat{X}=\mat{A}$ and $\mat{Y}=(\mat{A}\T\mat{A})\inv\mat{A}\T$; apply the cyclic property to push
$\mat{A}$ to the right, and watch the expression collapse to the trace of a $2\times 2$ identity matrix.)
\end{exercise}

\begin{exercise}\label{ch18:ex-complement}
Let $\mat{P}$ project onto a subspace $S$. Prove that $\mat{Q}=\mat{I}-\mat{P}$ is also a projection
matrix (symmetric and idempotent), and argue geometrically that $\mat{Q}$ projects onto the
orthogonal complement $S^\perp$. What is $\mat{P}\mat{Q}$, and why does the answer make geometric
sense?
\end{exercise}

\begin{exercise}\label{ch18:ex-closest-num}
For the line through $\vec{a}=\colvec{2\\1}$ and the point $\vec{b}=\colvec{1\\4}$ (from
Example~\ref{ch18:line-diagonal}), pick any \emph{other} point on the line, say $\vec{y}=2\vec{a}$,
and verify numerically that $\norm{\vec{b}-\vec{y}} > \norm{\vec{b}-\vec{p}}$. Then confirm the
Pythagorean identity $\norm{\vec{b}-\vec{y}}^2 = \norm{\vec{e}}^2 + \norm{\vec{p}-\vec{y}}^2$ with
your numbers.
\end{exercise}

\begin{exercise}\label{ch18:ex-eigen}
\textbf{(Preview of eigenvalues.)} A scalar $\lambda$ is called an \emph{eigenvalue} of a square
matrix $\mat{P}$ if there is a nonzero vector $\vec{v}$ that $\mat{P}$ merely scales without rotating:
$\mat{P}\vec{v} = \lambda\vec{v}$. Using only the identity $\mat{P}^2=\mat{P}$, prove that every
such scale factor must satisfy $\lambda^2=\lambda$, hence $\lambda\in\set{0,1}$. (Hint: apply $\mat{P}$
to both sides of $\mat{P}\vec{v}=\lambda\vec{v}$.) Which vectors correspond to $\lambda=1$, and which
to $\lambda=0$?
\end{exercise}

\begin{exercise}\label{ch18:ex-oblique}
\textbf{(Challenge: why symmetry matters.)} Consider $\mat{P}=\begin{bmatrix}1 & 1\\ 0 & 0\end{bmatrix}$.
Verify that $\mat{P}^2=\mat{P}$, so $\mat{P}$ is idempotent, but $\mat{P}\T\neq\mat{P}$. Show that
$\mat{P}$ maps every vector onto the $x$-axis, yet does \emph{not} send $\vec{b}=\colvec{0\\1}$ to
its \emph{closest} point on the $x$-axis (which is the origin). Conclude that idempotence alone
gives a projection, but only \emph{symmetry} makes it the orthogonal, closest-point projection.
\end{exercise}

\begin{exercise}\label{ch18:ex-orthonormal}
\textbf{(Challenge.)} Suppose the columns of $\mat{A}$ are pairwise \emph{orthonormal}:
$\mat{A}\T\mat{A}=\mat{I}$. Show that the projection matrix simplifies to
$\mat{P}=\mat{A}\mat{A}\T$, and that $\hat{\vec{x}}=\mat{A}\T\vec{b}$ -- each coefficient is just a
single dot product $\vec{a}_i\T\vec{b}$, with no matrix inverse anywhere. Explain why this is the
computational reason the next-but-one chapter works so hard to build orthonormal bases via
Gram--Schmidt.
\end{exercise}
